% Nouveau EGA V (cote 25): archive cote 25, site page 122 (PDF page 123).
% typescript. Draft compared: none.
% AI-transcribed by Claude Opus 5.5 (claude-opus-5-5, adaptive thinking, effort high), Message Batches API, batch ega25b (msgbatch_01A7n5YuXubUPNK1zsepi8RG), 2026-11-13; checked against both supplied scan parts; see CHECK_DECISIONS.tsv.
\scanpage{25/122}
\mspage{25}

systêmes linéaires de diviseurs sans composante fixe  sur X paramétrés
par $\overset{\hand{\vee}}{P}$ , et l'ensemble $\hand{\unsure{\psi|\emptyset}}$  des applications rationelles de X dans P, \hand{qui}
\var{\inserted{sur $\bar{k}$ ne se factorisent par aucun hyperplan de $\unsure{\overline{P}|P}_{\bar{k}}$.}}{sur $\bar{k}$ ne se factorisent par aucun hyperplan de $\unsure{\overline{P}|P}_{\bar{k}}$.}{insertion}{ega25_25_122-v0001}


En effet, \var{\inserted[typed]{lorsque les anneaux locaux de X sont factoriels.}}{lorsque les anneaux locaux de X sont factoriels.}{insertion}{ega25_25_122-v0002} \var{\deleted{\var{\xout{X étant régulier, \unsure{il|X} en est de même de \unsure{$X\times_S P$} qui est lisse}}{}{typewriter strike-out}{ega25_25_122-v0004}
\var{\xout{sur X, donc d'après un résultat de Auslander-Buchsbaum (qui sera repris}}{}{typewriter strike-out}{ega25_25_122-v0005}
\var{\xout{au Chap III) \illegible{1} de les anneaux locaux de \unsure{$X\times_S P$} sont factoriels,}}{}{typewriter strike-out}{ega25_25_122-v0006}
et comme codim(X-U,X) \quad 2}}{}{deletion}{ega25_25_122-v0003} , il s'ensuit que tout Module inversible sur
U peut se prolonger en un Module inversible sur X, donc la condition
envisagée dans b) est automatiquement satisfaite. D'autre part, d'après
\var{\deleted{\unsure{un Mémoire} \illegible{1} \unsure{dans un} n° \unsure{ultérieur}}}{}{deletion}{ega25_25_122-v0007}
Auslander-Buchsbaum (\var{\deleted{qui figurera un jour au Chap III}}{}{deletion}{ega25_25_122-v0008}\var{\inserted{cf §20}}{cf §20}{insertion}{ega25_25_122-v0009}) un annea\hand{/}u local
régulier est factoriel.

\var{\margin{\unsure{$\times\quad\times\quad\nearrow$}}}{\placedmargin{\unsure{$\times\quad\times\quad\nearrow$}}}{marginal note placed}{ega25_25_122-v0010}La proposition 2  b)  nous montre que myennant une hypothèse
\var{\inserted[typed]{(savoir $D\quad\hand{\unsure{\phi|\emptyset}}$)}}{(savoir $D\quad\hand{\unsure{\phi|\emptyset}}$)}{insertion}{ega25_25_122-v0011}
supplémentaire sur un\var{\xout{e famille}}{}{typewriter strike-out}{ega25_25_122-v0012} système linéaire de diviseurs D, on peut
lui associer canoniquement un Module inversible $\underline{L}$ sur X, prolongement
\var{\inserted[typed]{où U=X-Z.}}{où U=X-Z.}{insertion}{ega25_25_122-v0013}
canonique de $\underline{L}^{x}_{U} = f\ (\underline{O}_P(1))$ ,  On peut en donner une autre interprétation,
indépendante de l'hypothèse  $D\quad\hand{\unsure{\phi|\emptyset}}$ , de la façon suivante; par constructio

\var{\xout{l'image inverse de $\underline{L}$x}}{}{typewriter strike-out}{ega25_25_122-v0014} ou presque $\underline{L}^{x}_{U}\otimes_S\underline{O}_P(1)$ \var{\xout{n'}}{}{typewriter strike-out}{ega25_25_122-v0015}est \var{\xout{autre que}}{}{typewriter strike-out}{ega25_25_122-v0016} \var{\inserted[typed]{can isomorphe à}}{can isomorphe à}{insertion}{ega25_25_122-v0017} la restric-%
tion $\underline{N}'$ à $U\mathrm{x}_S\overset{\hand{\vee}}{P}$  du  Module inversible $\underline{N}$  sur $X\mathrm{x}_S\overset{\hand{\vee}}{P}$  défini par D. En vertu
du \var{\xout{n°1}}{}{typewriter strike-out}{ega25_25_122-v0018} n°1, cela détermine \var{\xout{uniquement $\underline{N}$, et}}{}{typewriter strike-out}{ega25_25_122-v0019} $\underline{L}_U$  à isom unique près;
et  (comme U est dense dans X)  il existe  un Module inversible $\underline{L}$ sur X,
déterminé à isom unique près (do\unsure{n|ç}c prolongeant $\underline{L}_U$ ) tel que  l'on ait
un isomorphisme  $\underline{L}\otimes_{\underline{O}_S}\underline{O}_{\overset{\hand{\unsure{\vee|/}}}{P}}\ (1) \quad \underline{N}$ ;  $\underline{L}$  est donné par exemple par
\[ \underline{L} \quad pr_1\ (\underline{N}) \quad ; \]
on l'appelle  le Module inversible sur X associé au systême lin de divi-%
suurs D.

\mspage{122}
